% Original: PDF strana 11, gornji slajd; drugi deo.
\slajdid{03-s021b}
\begin{frame}[t,label=03-s021b]{Crtanje simbola kola srednjeg stepena integracije}
\setlength{\parskip}{0pt}
% element: 03-s021b:e04
Verovatno ste uočili da nema strogog standarda za crtanje simbola
kola srednjeg stepena integracije.
\par\vspace{3mm}
\begin{columns}[c,onlytextwidth]
\begin{column}{.36\textwidth}\centering
\Oznaka{Pravougaoni simbol komponente}
\par\vspace{1mm}
\begin{tikzpicture}[x=1mm,y=1mm,font=\fontsize{9}{11}\selectfont]
\draw[DEvod] (0,0) rectangle (28,32);
\node at (14,23) {KP};\node at (14,18) {4/2};
\foreach \lab/\yy in {I3/26,I2/21,I1/16,I0/11,EI/6}{
 \node[anchor=west,inner sep=1pt] at (0,\yy) {\lab};
 \draw[DEvod] (-3,\yy)--(0,\yy);
}
\foreach \lab/\yy in {GS/26,A1/20,A0/14,EO/6}{
 \node[anchor=east,inner sep=1pt] at (28,\yy) {\lab};
 \draw[DEvod] (28,\yy)--(31,\yy);
}
\end{tikzpicture}

\end{column}
\begin{column}{.60\textwidth}
\textcolor{DEisticanje}{\textbf{Uobičajeno je da svi ulazi budu na jednoj strani
simbola, a svi izlazi na drugoj.}}
\par\vspace{3mm}
\textbf{Ulaze i izlaze ne mešamo na istoj stranici simbola.}
\par\vspace{3mm}
Dosta često crtamo \textbf{ulaze sa leve}, a \textbf{izlaze sa desne strane}.
\end{column}
\end{columns}
\end{frame}
